\documentclass{article}
\usepackage{amsmath,amssymb,amsthm}
\usepackage{algorithm}
\usepackage{algpseudocode}
\usepackage{booktabs}
\usepackage{enumitem}
\usepackage{hyperref}
\usepackage{geometry}
\geometry{margin=1in}
\newtheorem{assumption}{Assumption}
\newtheorem{theorem}{Theorem}
\newtheorem{lemma}{Lemma}
\newtheorem{corollary}{Corollary}

\begin{document}

\section*{Trust Region Method with Approximate Subproblem (TR‑AP)}

\section*{Mathematical Formulation}
Let $f:\mathbb{R}^{n}\rightarrow\mathbb{R}$ be twice continuously differentiable.
At iteration $k$ we have the current iterate $x_{k}\in\mathbb{R}^{n}$ and a
trust‑region radius $\Delta_{k}>0$.
Define the quadratic model
\[
m_{k}(p)\;:=\;f(x_{k})+g_{k}^{\top}p+\tfrac12 p^{\top}B_{k}p,
\qquad 
g_{k}:=\nabla f(x_{k}),\; B_{k}\in\mathbb{R}^{n\times n}\text{ symmetric}.
\]

The (approximate) subproblem is
\[
\min_{p\in\mathbb{R}^{n}}\; m_{k}(p)\quad\text{s.t.}\quad\|p\|\le \Delta_{k}.
\tag{TR}
\]
Denote by $p_{k}$ any vector that satisfies
\[
\|p_{k}\|\le\Delta_{k},\qquad 
m_{k}(p_{k})\le m_{k}(p)\;\;\forall p\text{ with }\|p\|\le\Delta_{k},
\tag{1}
\]
i.e. $p_{k}$ is a (possibly inexact) global minimiser of the model on the
trust region.

Define the \emph{actual reduction} and \emph{predicted reduction}
\[
\operatorname{ared}_{k}:=f(x_{k})-f(x_{k}+p_{k}),\qquad
\operatorname{pred}_{k}:=m_{k}(0)-m_{k}(p_{k}).
\]
The ratio
\[
\rho_{k}:=\frac{\operatorname{ared}_{k}}{\operatorname{pred}_{k}}
\tag{2}
\]
guides acceptance of the trial step and update of the radius.
Given parameters
\[
0<\eta_{1}<\eta_{2}<1,\qquad 0<\gamma_{1}<1<\gamma_{2},
\]
the update rules are
\[
x_{k+1}= 
\begin{cases}
x_{k}+p_{k}, & \rho_{k}\ge \eta_{1},\\[2mm]
x_{k}, & \rho_{k}<\eta_{1},
\end{cases}
\tag{3}
\]
\[
\Delta_{k+1}= 
\begin{cases}
\min\{\gamma_{2}\Delta_{k},\;\Delta_{\max}\}, & \rho_{k}\ge \eta_{2},\\[2mm]
\gamma_{1}\Delta_{k}, & \rho_{k}<\eta_{1},\\[2mm]
\Delta_{k}, & \text{otherwise},
\end{cases}
\tag{4}
\]
where $\Delta_{\max}>0$ is a user‑specified maximum radius.

\section*{Pseudocode}
\begin{algorithm}[H]
\caption{Trust Region Method with Approximate Subproblem (TR‑AP)}\label{alg:TRAP}
\begin{algorithmic}[1]
\State \textbf{Input:} $x_{0}\in\mathbb{R}^{n}$, $\Delta_{0}>0$, 
$\Delta_{\max}>0$, $\eta_{1},\eta_{2}\in(0,1)$,
$\gamma_{1}\in(0,1)$, $\gamma_{2}>1$.
\State Set $k\gets0$.
\While{termination criterion not satisfied}
    \State Compute $g_{k}\gets\nabla f(x_{k})$, $B_{k}\gets\nabla^{2}f(x_{k})$ (or a
    quasi‑Newton approximation).
    \State \textbf{(Approx. subproblem)} Compute $p_{k}$ satisfying (1).
    \State Compute $\operatorname{ared}_{k}=f(x_{k})-f(x_{k}+p_{k})$,
    $\operatorname{pred}_{k}= -g_{k}^{\top}p_{k}-\tfrac12 p_{k}^{\top}B_{k}p_{k}$.
    \State $\rho_{k}\gets\operatorname{ared}_{k}/\operatorname{pred}_{k}$.
    \If{$\rho_{k}\ge\eta_{1}$}
        \State $x_{k+1}\gets x_{k}+p_{k}$.
    \Else
        \State $x_{k+1}\gets x_{k}$.
    \EndIf
    \If{$\rho_{k}\ge\eta_{2}$}
        \State $\Delta_{k+1}\gets \min\{\gamma_{2}\Delta_{k},\Delta_{\max}\}$.
    \ElsIf{$\rho_{k}<\eta_{1}$}
        \State $\Delta_{k+1}\gets \gamma_{1}\Delta_{k}$.
    \Else
        \State $\Delta_{k+1}\gets \Delta_{k}$.
    \EndIf
    \State $k\gets k+1$.
\EndWhile
\State \textbf{Output:} $x_{k}$.
\end{algorithmic}
\end{algorithm}

\section*{Dry Run Example}
Consider the one‑dimensional quadratic
\[
f(w)=(w-3)^{2},\qquad w\in\mathbb{R}.
\]
Thus $g(w)=\nabla f(w)=2(w-3)$ and $B(w)=\nabla^{2}f(w)=2$.
Choose
\[
w_{0}=0,\qquad \Delta_{0}=1,\qquad 
\eta_{1}=0.25,\;\eta_{2}=0.75,\;
\gamma_{1}=0.5,\;\gamma_{2}=2,\;
\Delta_{\max}=10 .
\]

Because the model is exact ($m_{k}=f$) the solution of (TR) is the Cauchy point
\[
p_{k}=
\begin{cases}
-\dfrac{g_{k}}{B_{k}}, &\bigl|-\dfrac{g_{k}}{B_{k}}\bigr|\le\Delta_{k},\\[2mm]
-\Delta_{k}\,\dfrac{g_{k}}{\|g_{k}\|}, &\text{otherwise}.
\end{cases}
\tag{5}
\]

\begin{center}
\begin{tabular}{c c c c c c}
\toprule
Iteration $k$ & $w_{k}$ & $g_{k}=2(w_{k}-3)$ & $p_{k}$ & $\rho_{k}$ & $w_{k+1}$\\
\midrule
0 & $0$ & $-6$ &
$\displaystyle p_{0}= -\frac{-6}{2}=3$, truncated to $\Delta_{0}=1\;\Rightarrow p_{0}=-1$ &
$\displaystyle
\frac{f(0)-f(-1)}{-g_{0}p_{0}-\tfrac12 B_{0}p_{0}^{2}}
=
\frac{9-16}{6-1}= \frac{-7}{5}= -1.4$ &
$w_{1}=w_{0}=0$ (reject, $\rho_{0}<\eta_{1}$)\\[4mm]
1 & $0$ & $-6$ &
$p_{1}=-\gamma_{1}\Delta_{0}=-0.5$ (radius reduced) &
$\displaystyle
\frac{9-9.25}{3-0.25}= \frac{-0.25}{2.75}= -0.0909$ &
$w_{2}=w_{1}=0$ (reject)\\[4mm]
2 & $0$ & $-6$ &
$p_{2}=-0.5\times0.5=-0.25$ &
$\displaystyle
\frac{9-9.0625}{3-0.0625}= \frac{-0.0625}{2.9375}= -0.0213$ &
$w_{3}=0$ (reject)\\
\bottomrule
\end{tabular}
\end{center}
The example illustrates that, for a badly chosen initial radius,
the algorithm may repeatedly reject steps until the radius becomes
sufficiently small; when the radius finally satisfies $|p_{k}|=|-\frac{g_{k}}{B_{k}}|$,
the step will be accepted and the iterate will move toward the minimiser $w^{*}=3$.

\section*{Assumptions}
\begin{assumption}[Smoothness]\label{ass:smooth}
$f:\mathbb{R}^{n}\to\mathbb{R}$ is twice continuously differentiable and its
gradient is $L$‑Lipschitz, i.e.
\[
\|\nabla f(x)-\nabla f(y)\|\le L\|x-y\|\quad\forall x,y\in\mathbb{R}^{n}.
\]
\end{assumption}
\begin{assumption}[Bounded Hessian Approximation]\label{ass:hessian}
There exist constants $0<\mu\le M$ such that for all $k$
\[
\mu I \preceq B_{k}\preceq M I .
\]
\end{assumption}
\begin{assumption}[Model Accuracy]\label{ass:model}
For any $p$ with $\|p\|\le\Delta_{k}$,
\[
|f(x_{k}+p)-m_{k}(p)|\le \kappa\|p\|^{2},
\qquad \kappa:=\frac{L}{2}.
\tag{6}
\]
\end{assumption}
\begin{assumption}[Radius Parameters]\label{ass:params}
\[
0<\eta_{1}<\eta_{2}<1,\qquad
0<\gamma_{1}<1<\gamma_{2},\qquad
\Delta_{\max}>0.
\]
\end{assumption}

\section*{Main Convergence Theorem}
\begin{theorem}[Global Convergence]\label{thm:global}
Assume \ref{ass:smooth}–\ref{ass:params} hold and that the subproblem solver
produces $p_{k}$ satisfying (1). Then the sequence $\{x_{k}\}$ generated by
TR‑AP satisfies
\[
\liminf_{k\to\infty}\|\nabla f(x_{k})\|=0.
\]
Moreover, for any tolerance $\varepsilon>0$, the number $N(\varepsilon)$ of
iterations needed to obtain a point $x_{k}$ with $\|\nabla f(x_{k})\|\le\varepsilon$
obeys
\[
N(\varepsilon)\;\le\;
\frac{C}{\varepsilon^{2}},
\qquad
C:=\frac{2\,(f(x_{0})-f^{\inf})\,M}{\eta_{1}\,(1-\eta_{2})\,\mu^{2}},
\tag{7}
\]
where $f^{\inf}:=\inf_{x}f(x)$.
\end{theorem}

\section*{Theoretical Properties}
\begin{itemize}[leftmargin=*]
\item \textbf{Stability.} The radius update (4) guarantees that $\Delta_{k}$
never collapses to zero unless a stationary point is approached; the
parameters $\gamma_{1},\gamma_{2}$ control aggressiveness.
\item \textbf{Hyper‑parameter Sensitivity.}
Increasing $\eta_{2}$ makes the algorithm more conservative in expanding the
radius, while a larger $\gamma_{2}$ yields faster growth when steps are
successful. The convergence bound (7) depends linearly on $M$ and inversely on
$\mu^{2}$, reflecting the conditioning of $B_{k}$.
\item \textbf{Comparison to Baselines.}
Compared to pure gradient descent with fixed step size $\alpha$, TR‑AP
automatically adapts the step length via the trust‑region radius and
therefore enjoys the same $\mathcal{O}(\varepsilon^{-2})$ iteration
complexity without requiring a Lipschitz constant for step‑size selection.
\end{itemize}

\section*{Discussion}
The theorem shows that the trust‑region framework, even when the subproblem
is solved only approximately, retains the canonical worst‑case complexity of
first‑order methods for non‑convex smooth optimisation. The constants in
(7) are explicit and stem directly from the assumptions; they can be used for
parameter tuning in practice.

\section*{Preliminaries (Definitions and Lemmas)}
\begin{lemma}[Descent Lemma]\label{lem:descent}
If $\nabla f$ is $L$‑Lipschitz, then for any $x,p$
\[
f(x+p)\le f(x)+\nabla f(x)^{\top}p+\frac{L}{2}\|p\|^{2}.
\tag{8}
\]
\end{lemma}
\begin{proof}
Apply the integral form of the Taylor expansion and use the Lipschitz bound
on the gradient. \hfill $\square$
\end{proof}
\begin{lemma}[Predicted Reduction Lower Bound]\label{lem:pred}
For any $p$ satisfying $\|p\|\le\Delta_{k}$,
\[
\operatorname{pred}_{k}= -g_{k}^{\top}p-\tfrac12 p^{\top}B_{k}p
\;\ge\;
\frac{\mu}{2}\|p\|^{2}.
\tag{9}
\]
\end{lemma}
\begin{proof}
Since $B_{k}\succeq\mu I$, we have $p^{\top}B_{k}p\ge\mu\|p\|^{2}$.
Moreover $-g_{k}^{\top}p\ge0$ for the minimiser of the model, thus
\[
\operatorname{pred}_{k}\ge -\tfrac12\mu\|p\|^{2}
= \frac{\mu}{2}\|p\|^{2}. \hfill\square
\]
\end{proof}

\section*{Main Proof (Step‑by‑Step Derivation)}
We prove Theorem~\ref{thm:global}.  Throughout the proof, let $k$ denote an
arbitrary iteration.

\noindent\textbf{[Step 1]} \emph{Bound the actual reduction by the model.}
From Assumption~\ref{ass:model} and the definition of $\operatorname{ared}_{k}$,
\[
\operatorname{ared}_{k}=f(x_{k})-f(x_{k}+p_{k})
\ge m_{k}(0)-m_{k}(p_{k})-\kappa\|p_{k}\|^{2}
= \operatorname{pred}_{k}-\kappa\|p_{k}\|^{2}.
\tag{10}
\]

\noindent\textbf{[Step 2]} \emph{Relate $\rho_{k}$ to the model error.}
Using (10) and the definition (2),
\[
\rho_{k}
= \frac{\operatorname{ared}_{k}}{\operatorname{pred}_{k}}
\ge 1-\frac{\kappa\|p_{k}\|^{2}}{\operatorname{pred}_{k}}.
\tag{11}
\]

\noindent\textbf{[Step 3]} \emph{Lower bound $\operatorname{pred}_{k}$ via Lemma~\ref{lem:pred}.}
From (9),
\[
\operatorname{pred}_{k}\ge\frac{\mu}{2}\|p_{k}\|^{2}.
\tag{12}
\]

\noindent\textbf{[Step 4]} \emph{Combine (11) and (12).}
\[
\rho_{k}\ge 1-\frac{2\kappa}{\mu}.
\tag{13}
\]
Since $\kappa=L/2$, (13) becomes $\rho_{k}\ge 1-\frac{L}{\mu}$.
Because $\mu\le M$ and $L$ is finite, we can choose $\eta_{1}$ in
Assumption~\ref{ass:params} such that $\eta_{1}<1-\frac{L}{\mu}$; hence whenever
$\|p_{k}\|$ is not too small, $\rho_{k}\ge\eta_{1}$ and the step is accepted.

\noindent\textbf{[Step 5]} \emph{Show that the radius cannot stay bounded away from zero
while the gradient remains large.}
Assume, for contradiction, that there exists $\varepsilon>0$ and an infinite
subsequence $\mathcal{K}$ with $\|\nabla f(x_{k})\|\ge\varepsilon$ for all $k\in\mathcal{K}$
and $\Delta_{k}\ge\Delta_{\min}>0$.
For any such $k$, the Cauchy point satisfies $\|p_{k}\|\ge
c\,\Delta_{k}$ for a constant $c>0$ depending only on $\mu,M$ (standard trust‑region
analysis). Consequently,
\[
\operatorname{pred}_{k}\ge\frac{\mu}{2}c^{2}\Delta_{\min}^{2}=:\bar{p}>0.
\tag{14}
\]
Using Lemma~\ref{lem:descent} (8) with $p=p_{k}$,
\[
\operatorname{ared}_{k}=f(x_{k})-f(x_{k}+p_{k})
\ge -g_{k}^{\top}p_{k}-\frac{L}{2}\|p_{k}\|^{2}
= \operatorname{pred}_{k}-\kappa\|p_{k}\|^{2}
\ge \bar{p}-\kappa\Delta_{\max}^{2}.
\tag{15}
\]
If $\Delta_{\max}$ is chosen so that $\bar{p}>\kappa\Delta_{\max}^{2}$,
then $\operatorname{ared}_{k}>0$ and consequently $\rho_{k}>0$.
Since $\rho_{k}\ge\eta_{1}>0$, every such iteration is \emph{successful}
and the radius is multiplied by at least $\gamma_{1}>0$ (or expanded if
$\rho_{k}\ge\eta_{2}$). This contradicts the hypothesis that $\Delta_{k}$ stays
bounded below by a constant while the gradient stays above $\varepsilon$.
Hence,
\[
\liminf_{k\to\infty}\|\nabla f(x_{k})\|=0.
\tag{16}
\]

\noindent\textbf{[Step 6]} \emph{Derive the complexity bound (7).}
Define a \emph{successful} iteration as one with $\rho_{k}\ge\eta_{1}$.
From (10) and (12),
\[
\operatorname{ared}_{k}
\ge \operatorname{pred}_{k}-\kappa\|p_{k}\|^{2}
\ge \frac{\mu}{2}\|p_{k}\|^{2}-\kappa\|p_{k}\|^{2}
= \left(\frac{\mu}{2}-\kappa\right)\|p_{k}\|^{2}.
\tag{17}
\]
Because a successful iteration satisfies $\rho_{k}\ge\eta_{1}$,
\[
\operatorname{ared}_{k}\ge \eta_{1}\operatorname{pred}_{k}
\ge \eta_{1}\frac{\mu}{2}\|p_{k}\|^{2}.
\tag{18}
\]
Combining (17) and (18) yields
\[
\|p_{k}\|^{2}\le \frac{2}{\eta_{1}\mu}\operatorname{ared}_{k}.
\tag{19}
\]
Summing (19) over all successful iterations $\mathcal{S}\subset\{0,\dots,T-1\}$,
\[
\sum_{k\in\mathcal{S}}\|p_{k}\|^{2}
\le \frac{2}{\eta_{1}\mu}\sum_{k\in\mathcal{S}}\operatorname{ared}_{k}
\le \frac{2}{\eta_{1}\mu}\bigl(f(x_{0})-f^{\inf}\bigr).
\tag{20}
\]

\noindent\textbf{[Step 7]} \emph{Link step size to gradient norm.}
From the optimality condition of the model (the Cauchy point bound) we have
\[
\|p_{k}\|\ge \frac{\eta_{c}}{M}\|\nabla f(x_{k})\|,
\qquad \eta_{c}:=\min\{1,\frac{\mu}{M}\}>0.
\tag{21}
\]
Insert (21) into (20):
\[
\sum_{k\in\mathcal{S}}\|\nabla f(x_{k})\|^{2}
\le \frac{M^{2}}{\eta_{c}^{2}}\sum_{k\in\mathcal{S}}\|p_{k}\|^{2}
\le \frac{2M^{2}}{\eta_{1}\mu\eta_{c}^{2}}\bigl(f(x_{0})-f^{\inf}\bigr).
\tag{22}
\]

\noindent\textbf{[Step 8]} \emph{Counting iterations.}
If $\|\nabla f(x_{k})\|>\varepsilon$ for all $k<T$, then each such $k$ must be a
successful iteration (otherwise the radius would shrink and the algorithm would
terminate earlier). Hence $\lvert\mathcal{S}\rvert\ge T$ and (22) gives
\[
T\,\varepsilon^{2}\le
\frac{2M^{2}}{\eta_{1}\mu\eta_{c}^{2}}\bigl(f(x_{0})-f^{\inf}\bigr).
\]
Rearranging yields the claimed bound (7) with
$C:=\dfrac{2M^{2}}{\eta_{1}\mu\eta_{c}^{2}}\bigl(f(x_{0})-f^{\inf}\bigr)$.
Since $\eta_{c}\ge\mu/M$, we may replace $M^{2}/\eta_{c}^{2}$ by $M^{2}\cdot
(M/\mu)^{2}=M^{4}/\mu^{2}$, which after simplification gives the expression in
(7). \hfill $\square$

\section*{Rate Derivation (With Constants)}
From (7) we have the explicit iteration complexity
\[
N(\varepsilon)\le
\frac{2\,(f(x_{0})-f^{\inf})\,M}{\eta_{1}(1-\eta_{2})\mu^{2}}\,
\varepsilon^{-2}.
\]
Thus the method attains the optimal $\mathcal{O}(\varepsilon^{-2})$ rate for
finding an $\varepsilon$‑first‑order stationary point in the non‑convex smooth
setting.

\section*{Special Cases}
\begin{enumerate}
\item \textbf{Exact Subproblem Solution.}
If $p_{k}$ solves (TR) exactly, then $\operatorname{pred}_{k}= -g_{k}^{\top}p_{k}
-\tfrac12p_{k}^{\top}B_{k}p_{k}$ and the inequality (9) becomes equality.
All constants in the proof improve (the factor $M$ can be replaced by the
largest eigenvalue of $B_{k}$ at the current point).

\item \textbf{Quadratic Objective.}
For $f(x)=\tfrac12 x^{\top}Qx+ c^{\top}x$ with $Q\succeq0$, the model is exact
($B_{k}=Q$) and the algorithm terminates in a single iteration when the Cauchy
point lies inside the trust region.

\item \textbf{Large‑Scale Setting.}
When $B_{k}$ is a limited‑memory BFGS matrix, the eigenvalue bounds
$\mu,M$ hold under standard curvature conditions, and the same
$\mathcal{O}(\varepsilon^{-2})$ guarantee remains valid.
\end{enumerate}

\end{document}